% Standalone wrapper for the population-transfer report section.
% Build: pdflatex main && bibtex main && pdflatex main && pdflatex main
\documentclass[11pt,a4paper]{article}
\usepackage[T1]{fontenc}
\usepackage[utf8]{inputenc}
\usepackage{lmodern}
\usepackage[margin=1in]{geometry}
\usepackage{amsmath,amssymb,amsthm}
\usepackage{booktabs}
\usepackage{graphicx}
\usepackage{microtype}
\usepackage[authoryear,round]{natbib}
\usepackage[font=small,labelfont=bf]{caption}
\usepackage[dvipsnames]{xcolor}
\usepackage[colorlinks=true,linkcolor=blue!50!black,citecolor=blue!50!black,urlcolor=blue!50!black]{hyperref}
\setlength{\parskip}{0.4em}
\setlength{\parindent}{0pt}

\title{\textbf{Enforcing Constraints for Policy Transfer Across Agent Populations and Spatial Scales}\\[2pt]
\large Individual contribution: population transfer ($N\to N{+}k$) with shared, independent and attention policies}

\begin{document}
\maketitle

\begin{abstract}
\noindent
A decentralised policy trained for a team of $N$ agents is often reused when the team grows. Task performance may
survive this change while the \emph{constraints} the policy was meant to respect do not. We study this in a controlled
cooperative-navigation world with a collision-avoidance constraint. We train three policy families on three agents
(independent per-agent networks, one shared network, and a shared attention network over neighbours) and test them
without retraining on three to six agents, measuring minimum inter-agent separation, collision rate and goal-reaching
success. Task success transfers almost perfectly, but constraint satisfaction does not: the fraction of episodes with a
collision roughly doubles when a fourth agent is added for the shared and independent policies. The attention policy
degrades far more slowly, and an ablation shows that this is due to the architecture rather than only to seeing more
neighbours.
\end{abstract}

% ==============================================================================================
%  population_transfer_section.tex
%  Individual contribution (Kripalu Sonar): population transfer N -> N+k with shared, independent
%  and attention policies.
%
%  Use inside the team's main.tex with   % ==============================================================================================
%  population_transfer_section.tex
%  Individual contribution (Kripalu Sonar): population transfer N -> N+k with shared, independent
%  and attention policies.
%
%  Use inside the team's main.tex with   % ==============================================================================================
%  population_transfer_section.tex
%  Individual contribution (Kripalu Sonar): population transfer N -> N+k with shared, independent
%  and attention policies.
%
%  Use inside the team's main.tex with   \input{population_transfer_section}
%  Needs these packages (ICLR style already loads most of them): amsmath, amssymb, amsthm,
%  booktabs, graphicx, natbib.  No tikz is needed: the architecture diagram is a PDF.
%  Figures (folder figures/): architectures.pdf, transfer_metrics.png, training_curves.png.
%  Citation keys are defined in refs_population_transfer.bib.
%  All numbers in the tables are generated from runs/results.json by the code repository.
% ==============================================================================================
\makeatletter
\@ifundefined{proposition}{\newtheorem{proposition}{Proposition}}{}
\@ifundefined{remark}{\theoremstyle{remark}\newtheorem{remark}{Remark}\theoremstyle{plain}}{}
\makeatother

\section{Population Transfer with Shared, Independent and Attention Policies}
\label{sec:pt}

\subsection{Overview and relation to the proposal}
\label{sec:pt-overview}
Our proposal asks whether a decentralised policy that satisfies a constraint on one configuration of agents
keeps satisfying it after the configuration changes. This section studies the \emph{population} axis of that
question, $N\to N{+}k$ agents, in the simplest setting in which it can be asked cleanly: cooperative
navigation with a collision-avoidance constraint. The world is deliberately small, so that every number can be
reproduced on a CPU in minutes and so that the effect of the \emph{policy architecture} on constraint
preservation is not hidden behind other design choices. 

\paragraph{Contributions.}
\begin{enumerate}
  \item A transfer protocol that judges a policy by its \emph{constraint-related} metrics and not by reward
        alone (Section~\ref{sec:pt-setting}).
  \item A comparison of three policy families that differ in how they handle a change in the number of agents:
        independent per-agent networks, a single shared network, and a shared attention network over neighbours,
        plus an ablation that separates architecture from information (Section~\ref{sec:pt-policies}).
  \item Empirical findings on how much of the constraint survives transfer and which claims our data do and do not
        support (Sections~\ref{sec:pt-results}--\ref{sec:pt-limits}).
\end{enumerate}

\paragraph{Metrics of the proposal.}
The proposal measures transfer through consensus error, minimum inter-agent separation and constraint-violation
rate. Our task is goal reaching and not consensus, so there is no consensus error here. We measure the other two
quantities, and we add goal-reaching success as the measure of task performance.

\subsection{Related work}
\label{sec:pt-related}
\paragraph{Independent learners and parameter sharing.}
Independent learners train one policy per agent \citep{tan1993}. Parameter sharing uses one network for all agents
and is the standard way to obtain a policy that is defined for any team size \citep{gupta2017}. Both are strong
baselines in cooperative tasks when trained with PPO \citep{dewitt2020,yu2022}. Centralised training with
decentralised execution \citep{lowe2017,foerster2016} is the setting assumed in our proposal; our agents are
trained with local critics and executed in a fully decentralised way.

\paragraph{Set-based and attention architectures.}
Mean-field methods \citep{yang2018}, attention-based critics \citep{iqbal2019} and permutation-invariant set
encoders \citep{zaheer2017} treat the neighbours of an agent as a set and therefore apply to any number of
neighbours. Our attention policy follows this idea with a single attention head \citep{vaswani2017}.

\paragraph{Transfer across agent populations.}
Transfer in reinforcement learning is surveyed by \citet{taylor2009}. Our focus is how a \emph{collision
constraint} survives a change in the \emph{number of agents}.

\subsection{Problem setting}
\label{sec:pt-setting}
\paragraph{World.}
$N$ agents move in the unit square $[0,1]^2$. Agent $i$ has position $p_i^t$ and a goal $g_i$. There are five
actions $\mathcal A=\{\text{stay},\text{up},\text{down},\text{left},\text{right}\}$ with displacements
$\delta(a)\in\{0,\pm s e_1,\pm s e_2\}$, so $\|\delta(a)\|\le s$ with step size $s=0.05$. Positions are updated
simultaneously,
\begin{equation}
  p_i^{t+1}=\Pi_{[0,1]^2}\!\big(p_i^t+\delta(a_i^t)\big),
\end{equation}
where $\Pi$ projects back into the arena. An episode lasts $T=50$ steps. Starting positions and goals are
sampled uniformly, rejecting samples in which two starts or two goals are closer than $0.15$.

\paragraph{Constraint.}
Two agents \emph{collide} when $\|p_i^t-p_j^t\|<d_c$ with $d_c=0.1$. The constraint is
\begin{equation}
  \|p_i^t-p_j^t\|\ \ge\ d_c\qquad\text{for all } i\neq j \text{ and all } t.
\end{equation}

\paragraph{Observations.}
Every agent acts only on local information. With $r_{ij}=p_j-p_i$ and $\mathcal N_i^k$ the $k$ agents nearest to $i$,
\begin{equation}
  o_i^t=\big(p_i^t,\; g_i-p_i^t,\; (r_{ij}^t)_{j\in\mathcal N_i^k}\big)\in\mathbb R^{4+2k}.
\end{equation}
The shared and independent policies use $k=2$, so the observation length does not depend on $N$; this is what
allows a network trained with $N$ agents to be executed with $N{+}1$. The attention policy receives the
neighbours as a set whose size may vary (all $N{-}1$ other agents, or only the two nearest in the ablation).

\paragraph{Reward.}
Each agent receives
\begin{equation}
  r_i^t=10\big(\|g_i-p_i^{t-1}\|-\|g_i-p_i^t\|\big)\;-\;\lambda\,\mathbf{1}\!\big[\min_{j\neq i}\|p_i^t-p_j^t\|<d_c\big]\;+\;0.2\,\mathbf{1}\!\big[\|g_i-p_i^t\|<0.05\big],
\end{equation}
with collision penalty $\lambda=10$. The collision term is the only coupling between agents.

\paragraph{Transfer protocol.}
Each policy is trained on $N_{\text{train}}=3$ agents. It is then evaluated \emph{without any retraining} on
$N\in\{3,4,5,6\}$ agents in the same arena, so the team becomes denser as $N$ grows. As a reference we also train
each policy family directly on 4 agents. The difference between the transferred and the directly trained policy
at 4 agents is the \emph{transfer gap} (Section~\ref{sec:pt-results}).

\paragraph{Metrics.}
For each setting we run 500 evaluation episodes on the same random worlds for all policies (a paired comparison)
and use greedy (arg\,max) actions. We report means and standard deviations over three training seeds:
\begin{itemize}
  \item \textbf{Minimum separation:} $\min_{t,\,i\neq j}\|p_i^t-p_j^t\|$ within an episode, averaged over episodes
        (higher is safer; values below $d_c$ mean at least one collision).
  \item \textbf{Violation rate:} $\frac1T\sum_t\mathbf{1}[\exists\, i\neq j:\ \|p_i^t-p_j^t\|<d_c]$, the fraction of
        time steps with a collision.
  \item \textbf{Episodes with a collision:} the fraction of episodes with at least one collision.
  \item \textbf{Success rate:} the fraction of agents with $\|g_i-p_i^T\|<0.05$ at the final step
        (task performance).
\end{itemize}
A hand-coded \emph{greedy} controller that walks straight to its goal and ignores all other agents is included as a
``no avoidance'' reference. It shows how much of the loss of safety is due to the increased density alone.

\subsection{Policy architectures}
\label{sec:pt-policies}
\begin{figure}[t]
  \centering
  \includegraphics[width=0.78\linewidth]{figures/architectures.pdf}
  \caption{The three policy families, shown for a team trained on three agents that gets a fourth agent at test
  time (dashed). (a)~One network is shared by all agents. (b)~Each agent has its own network; the new agent has
  none and borrows a copy of agent~1's network. (c)~One shared network encodes the neighbours of an agent as a set
  and combines them with attention, so an extra neighbour only adds one element to the set.}
  \label{fig:pt-arch}
\end{figure}

\paragraph{Shared policy.}
One actor $\pi_\theta(a\mid o_i)$ and one critic $V_\phi(o_i)$, each a multilayer perceptron with two hidden layers
of 64 $\tanh$ units, are applied to every agent \citep{gupta2017}. An additional agent simply uses the same
network, so the policy is defined for every $N$ by construction.

\paragraph{Independent policies.}
Agent $i$ has its own actor $\pi_{\theta_i}$ and critic $V_{\phi_i}$ \citep{tan1993}. For $N>N_{\text{train}}$
there is no trained network for the additional agents, so an extension rule is needed. We let agent $i$ use the
network $\sigma(i)=((i-1)\bmod N_{\text{train}})+1$, i.e.\ the fourth agent receives a copy of agent~1's network.
This is a design choice and not a property of the method (Section~\ref{sec:pt-limits}).

\paragraph{Attention policy.}
With $r_{ij}=p_j-p_i$, the policy of agent $i$ computes
\begin{align}
  h_i &= \tanh\!\big(W_s[\,p_i;\ g_i-p_i\,]+b_s\big), &
  e_{ij} &= \tanh\!\big(W_n[\,r_{ij};\ \|r_{ij}\|\,]+b_n\big),\\
  \alpha_{ij} &= \operatorname*{softmax}_{j}\ \frac{(W_q h_i)^{\top}(W_k e_{ij})}{\sqrt d}, &
  c_i &= \sum_{j\neq i}\alpha_{ij}\,W_v e_{ij},
\end{align}
and outputs $\pi(\cdot\mid o_i)=\operatorname{softmax}\big(f_\theta([h_i;c_i])\big)$ with $d=64$ and $f_\theta$ a
two-layer MLP; the critic has the same structure. Because $c_i$ is a weighted average over the neighbours, it has
a fixed size and does not depend on the order of the neighbours: the network is permutation invariant and accepts
any number of neighbours \citep{zaheer2017}.

\paragraph{Ablation: attention on the two nearest neighbours.}
The attention policy sees all other agents, whereas the shared and independent policies see only the two nearest.
To separate architecture from information we also train the \emph{same} attention network on the two nearest
neighbours only. With three agents the two variants receive identical observations, so they coincide at $N=3$ by
construction; they can differ only for $N>3$.

\begin{table}[t]
  \centering\small
  \caption{Summary of the policies. ``Parameters'' counts actor and critic together for a team trained on three agents.}
  \label{tab:pt-arch}
  \begin{tabular}{l c l p{4.1cm} r}
    \toprule
    Policy & Networks & Neighbours seen & Extra agent at test time & Parameters \\
    \midrule
    Independent            & $N_{\text{train}}$ & 2 nearest & copy of the network of agent $\sigma(i)$ & 29,586 \\
    Shared                 & 1 & 2 nearest & same network & 9,862 \\
    Attention (all agents) & 1 & all other agents & same network & 38,278 \\
    Attention (2 nearest)  & 1 & 2 nearest & same network & 38,278 \\
    \bottomrule
  \end{tabular}
\end{table}

\subsection{Training}
\label{sec:pt-training}
All policies are trained with proximal policy optimisation \citep{schulman2017} and generalised advantage
estimation \citep{schulman2016}. Every agent uses only its own observation, both for the actor and for a local
critic, which corresponds to independent PPO \citep{dewitt2020}; only the actor is needed at execution time. One
iteration collects 50-step episodes in 64 parallel worlds and then updates the networks for four epochs. For the
independent policies the loss of every agent's network is computed on that agent's own experience. The settings are
listed in Table~\ref{tab:pt-hyper}.

\begin{table}[t]
  \centering\small
  \caption{Training settings (identical for all policy families).}
  \label{tab:pt-hyper}
  \begin{tabular}{ll}
    \toprule
    Quantity & Value \\
    \midrule
    Algorithm & PPO with GAE, local critics \\
    Optimiser & Adam, learning rate $10^{-3}$ \\
    Discount $\gamma$ / GAE $\lambda_{\text{GAE}}$ & 0.97 / 0.95 \\
    Clip range / entropy coefficient / value coefficient & 0.2 / 0.01 / 0.5 \\
    Gradient-norm clip & 0.5 \\
    Epochs per iteration / minibatch (world steps) & 4 / 512 \\
    Parallel worlds $\times$ steps per iteration & $64\times 50$ \\
    Iterations & 400 ($\approx$ 1.3 million world steps) \\
    Reward weights (progress / collision / goal) & 10 / 10 / 0.2 \\
    Training seeds & 0, 1, 2 \\
    \bottomrule
  \end{tabular}
\end{table}

\paragraph{The one tuned quantity.}
The collision penalty $\lambda$ was the only value we tuned. In a pilot with the shared policy, $\lambda=1$ and
$\lambda=3$ produced behaviour that was close to the greedy controller (agents had little incentive to avoid each
other), and $\lambda=10$ was the next value we tried and the first to give clear avoidance. We then used
$\lambda=10$ for all policies and all experiments. The pilot used the same evaluation worlds as the final
evaluation, which can make the absolute collision numbers slightly optimistic; it cannot favour one policy family,
since the value was fixed before the other families were trained.

\subsection{Experiments and results}
\label{sec:pt-results}
We train three seeds of every policy family on three agents (and, for the reference, three seeds of each of the
three main families on four agents) and evaluate on 3, 4, 5 and 6 agents.
All tables give the mean $\pm$ standard deviation over the three training seeds (population standard deviation).
With only three seeds we interpret large differences and treat differences of the order of the seed spread as
inconclusive.

\begin{figure}[t]
  \centering
  \includegraphics[width=\linewidth]{figures/transfer_metrics.png}
  \caption{Transfer from three to $N$ agents without retraining. Lines are means over three training seeds, faint dots
  are individual seeds, and the hollow diamonds are policies trained directly on four agents. \emph{Left:} minimum
  separation (collision below $0.1$). \emph{Middle:} fraction of time steps with a collision. \emph{Right:} fraction
  of agents on their goal at the end. The ``2 nearest'' attention variant (triangles) lies on top of the attention
  policy that sees all agents.}
  \label{fig:pt-transfer}
\end{figure}

\begin{table}[t]
  \centering\small
  \caption{Percentage of episodes with at least one collision (lower is safer).}
  \label{tab:pt-collision}
\begin{tabular}{l cccc}
\toprule
 & \multicolumn{4}{c}{Number of agents at test time} \\
\cmidrule(lr){2-5}
Policy (trained on 3 agents) & 3 (trained) & 4 & 5 & 6 \\
\midrule
Greedy (no avoidance) & 40.0 $\pm$ 0.0 & 66.6 $\pm$ 0.0 & 81.6 $\pm$ 0.0 & 92.8 $\pm$ 0.0 \\
\midrule
Independent & 20.7 $\pm$ 1.9 & 38.7 $\pm$ 1.4 & 58.4 $\pm$ 2.6 & 71.7 $\pm$ 2.2 \\
Shared & 16.3 $\pm$ 2.0 & 32.8 $\pm$ 3.9 & 49.5 $\pm$ 4.2 & 66.1 $\pm$ 3.4 \\
Attention (all agents) & 1.0 $\pm$ 0.2 & 3.0 $\pm$ 0.8 & 9.0 $\pm$ 3.1 & 15.3 $\pm$ 4.2 \\
Attention (2 nearest) & 1.0 $\pm$ 0.2 & 2.9 $\pm$ 1.1 & 8.3 $\pm$ 3.6 & 12.8 $\pm$ 5.7 \\
\midrule
\bottomrule
\end{tabular}
\end{table}

\begin{table}[t]
  \centering\small
  \caption{Minimum separation within an episode, averaged over episodes (higher is safer; the collision distance is
  $0.1$).}
  \label{tab:pt-minsep}
\begin{tabular}{l cccc}
\toprule
 & \multicolumn{4}{c}{Number of agents at test time} \\
\cmidrule(lr){2-5}
Policy (trained on 3 agents) & 3 (trained) & 4 & 5 & 6 \\
\midrule
Greedy (no avoidance) & 0.144 $\pm$ 0.000 & 0.089 $\pm$ 0.000 & 0.061 $\pm$ 0.000 & 0.045 $\pm$ 0.000 \\
\midrule
Independent & 0.192 $\pm$ 0.005 & 0.132 $\pm$ 0.003 & 0.095 $\pm$ 0.003 & 0.073 $\pm$ 0.004 \\
Shared & 0.201 $\pm$ 0.005 & 0.139 $\pm$ 0.005 & 0.104 $\pm$ 0.005 & 0.082 $\pm$ 0.004 \\
Attention (all agents) & 0.214 $\pm$ 0.008 & 0.172 $\pm$ 0.008 & 0.148 $\pm$ 0.006 & 0.131 $\pm$ 0.004 \\
Attention (2 nearest) & 0.214 $\pm$ 0.008 & 0.173 $\pm$ 0.008 & 0.150 $\pm$ 0.007 & 0.134 $\pm$ 0.007 \\
\midrule
\bottomrule
\end{tabular}
\end{table}

\begin{table}[t]
  \centering\small
  \caption{Percentage of time steps with a collision.}
  \label{tab:pt-violation}
\begin{tabular}{l cccc}
\toprule
 & \multicolumn{4}{c}{Number of agents at test time} \\
\cmidrule(lr){2-5}
Policy (trained on 3 agents) & 3 (trained) & 4 & 5 & 6 \\
\midrule
Greedy (no avoidance) & 2.80 $\pm$ 0.00 & 5.61 $\pm$ 0.00 & 8.76 $\pm$ 0.00 & 12.74 $\pm$ 0.00 \\
\midrule
Independent & 1.01 $\pm$ 0.11 & 2.07 $\pm$ 0.11 & 3.83 $\pm$ 0.25 & 5.64 $\pm$ 0.48 \\
Shared & 0.76 $\pm$ 0.11 & 1.71 $\pm$ 0.14 & 3.23 $\pm$ 0.09 & 4.96 $\pm$ 0.15 \\
Attention (all agents) & 0.04 $\pm$ 0.02 & 0.17 $\pm$ 0.09 & 0.42 $\pm$ 0.19 & 0.90 $\pm$ 0.31 \\
Attention (2 nearest) & 0.04 $\pm$ 0.02 & 0.16 $\pm$ 0.11 & 0.42 $\pm$ 0.23 & 0.95 $\pm$ 0.52 \\
\bottomrule
\end{tabular}
\end{table}

\begin{table}[t]
  \centering\small
  \caption{Success rate in percent (agents on their goal at the end of the episode).}
  \label{tab:pt-success}
\begin{tabular}{l cccc}
\toprule
 & \multicolumn{4}{c}{Number of agents at test time} \\
\cmidrule(lr){2-5}
Policy (trained on 3 agents) & 3 (trained) & 4 & 5 & 6 \\
\midrule
Greedy (no avoidance) & 100.0 $\pm$ 0.0 & 100.0 $\pm$ 0.0 & 100.0 $\pm$ 0.0 & 100.0 $\pm$ 0.0 \\
\midrule
Independent & 96.2 $\pm$ 0.6 & 96.3 $\pm$ 0.3 & 97.0 $\pm$ 0.2 & 97.1 $\pm$ 0.7 \\
Shared & 98.8 $\pm$ 0.5 & 98.9 $\pm$ 0.5 & 98.9 $\pm$ 0.6 & 98.8 $\pm$ 0.7 \\
Attention (all agents) & 97.4 $\pm$ 1.6 & 96.5 $\pm$ 1.6 & 94.9 $\pm$ 1.4 & 94.1 $\pm$ 1.3 \\
Attention (2 nearest) & 97.4 $\pm$ 1.6 & 96.2 $\pm$ 2.0 & 94.5 $\pm$ 1.9 & 92.9 $\pm$ 2.2 \\
\midrule
\bottomrule
\end{tabular}
\end{table}

\paragraph{Task success transfers, constraint satisfaction does not.}
Every learned policy reaches its goals in at least 92.9\,\% of the cases at every team size
(Table~\ref{tab:pt-success}), so the task itself transfers. The constraint does not. The share of episodes with a
collision (Table~\ref{tab:pt-collision}) is $16.3\,\%$ (shared) and $20.7\,\%$ (independent) with three agents and
roughly doubles when a fourth agent is added ($32.8\,\%$ and $38.7\,\%$); with six agents it is $66.1\,\%$ and
$71.7\,\%$. The mean minimum separation of the independent policies falls below the collision distance of $0.1$ at five
agents and that of the shared policy at six agents (Table~\ref{tab:pt-minsep}). A policy that looks successful by
its task reward can therefore already have lost much of its safety.

\paragraph{Part of the loss is density.}
The greedy controller, which has no avoidance at all, goes from $40.0\,\%$ to $92.8\,\%$ of episodes with a collision
between three and six agents, so a denser team is harder for any controller. The shared and independent policies
stay below the greedy controller at every size, but their advantage over it stays between $19$ and $34$ percentage points
and is smaller with six agents than with four: the increase from three to six agents is $49.8$ and $50.9$ points for
the shared and independent policies, against $52.8$ for the greedy controller. The advantage of the attention policy
over the greedy controller, in contrast, keeps growing, from $39$ to $78$ points.

\paragraph{The attention policy degrades much more slowly, and the cause is the architecture.}
The attention policy has $1.0\,\%$ of episodes with a collision at three agents and $15.3\,\%$ at six (an increase of
$14.3$ points), with a mean minimum separation of $0.131$ at six agents, still above $d_c$.
One might suspect that this is only because it sees all agents, whereas the other two see the two nearest. The
ablation rules this out for our setting: the same network fed only the two nearest neighbours reaches $12.8\,\%$ at
six agents (versus $15.3\,\%$; the difference is smaller than the seed spread of $4$--$6$ points) and is
indistinguishable at four and five agents (Table~\ref{tab:pt-collision}, Figure~\ref{fig:pt-transfer}). Seeing more
neighbours gave no measurable benefit. The comparison between the attention-on-two-nearest policy and the shared
policy differs in the architecture only, but ``architecture'' here is a bundle: set-structure with attention weights, an
explicit distance input $\|r_{ij}\|$, separate encoders, and about four times more parameters (Table~\ref{tab:pt-arch}).
We did not ablate these factors individually.

\paragraph{Shared versus independent.}
The shared policy is better than the independent policies in every column: $4.4$ to $8.9$ points fewer episodes with a
collision, a larger minimum separation and a higher success rate ($98.8\,\%$ versus $96.2\,\%$ at three agents). The
collision differences are of the same order as the seed spread ($1.4$--$4.2$ points), so we do not claim them to be
significant; the success difference is larger than the spread. Training curves
(Figure~\ref{fig:pt-train}) show that the shared policy learns faster than the independent ones, which is expected
since one network learns from the experience of all agents, and that the attention policy learns fastest.

\paragraph{Transfer gap.}
Table~\ref{tab:pt-gap} compares policies transferred from three to four agents with policies trained directly on four
agents. For the shared and independent policies the gap is large: $+17.7$ and $+12.9$ percentage points of episodes with a
collision. Directly trained policies are safer but succeed slightly less often, i.e.\ they trade some task success for
safety. For the attention policy there is no measurable gap ($-5.3$ points with a standard deviation of $8.7$); the
directly trained attention policies have $20.6\,\%$, $2.6\,\%$ and $1.8\,\%$ in the three seeds, so the poor average
comes from a single seed, and the transferred policy ($3.0\,\%$) is as safe as the two good ones. The reference runs use
the same training budget as the others and are therefore not necessarily converged.

\begin{table}[t]
  \centering\small
  \caption{Transfer gap at four agents: policies trained on three agents and transferred, versus policies trained
  directly on four agents. A positive gap means the transferred policy has more collisions (percentage points).}
  \label{tab:pt-gap}
\begin{tabular}{l cc c cc}
\toprule
 & \multicolumn{3}{c}{Episodes with a collision (\%)} & \multicolumn{2}{c}{Success rate (\%)} \\
\cmidrule(lr){2-4}\cmidrule(lr){5-6}
Policy & Transferred ($3\!\to\!4$) & Trained on 4 & Gap (pp) & Transferred & Trained on 4 \\
\midrule
Independent & 38.7 $\pm$ 1.4 & 25.8 $\pm$ 1.1 & +12.9 & 96.3 $\pm$ 0.3 & 92.0 $\pm$ 0.5 \\
Shared & 32.8 $\pm$ 3.9 & 15.1 $\pm$ 6.3 & +17.7 & 98.9 $\pm$ 0.5 & 96.9 $\pm$ 1.3 \\
Attention (all agents) & 3.0 $\pm$ 0.8 & 8.3 $\pm$ 8.7 & -5.3 & 96.5 $\pm$ 1.6 & 93.7 $\pm$ 4.0 \\
\bottomrule
\end{tabular}
\end{table}

\begin{figure}[t]
  \centering
  \includegraphics[width=0.92\linewidth]{figures/training_curves.png}
  \caption{Training curves on three agents (mean over three seeds, band = min to max; moving average of 10
  iterations). Left: training return. Right: fraction of time steps with a collision during training.}
  \label{fig:pt-train}
\end{figure}

\subsection{Discussion}
\label{sec:pt-discussion}
\paragraph{What this says about constraint-preserving transfer.}
Transfer of the task and transfer of the constraint are different things. The task survived a change of population
almost perfectly, while the constraint did not, and the loss was not predicted by task performance. For the
proposal this means that preservation of a constraint cannot be assumed from successful transfer; it has to be
measured directly. In our experiments, the set-based attention architecture preserved much more of the constraint
than the concatenated-observation policies, and its transfer gap at four agents was not measurable.

\paragraph{Information.}
When the team outgrows a $k$-nearest observation, each agent no longer observes every other agent. Our ablation
shows that, for the attention policy, observing all neighbours rather than only the two nearest did not improve the
measured collision outcomes. In this setting, the transfer gap is therefore not explained by this information
difference alone.

\paragraph{Practical reading.}
If a team may grow, the set-based attention policy is the most robust of the learned policies we tested: it has the
lowest collision rate while retaining high goal-reaching success across all evaluated team sizes.

\subsection{Limitations and threats to validity}
\label{sec:pt-limits}
\begin{itemize}
  \item \textbf{A small world.} The dynamics are discrete and noise free, there is no communication, and the arena is
        fixed, so the team becomes denser as it grows. Conclusions about larger teams, continuous dynamics or noise are
        not supported.
  \item \textbf{Three seeds.} We report means and standard deviations over three seeds without confidence intervals;
        only large effects (for example attention versus the other policy families) are reliable. The comparison between
        shared and independent policies is suggestive only.
  \item \textbf{Density and policy are confounded in $N$.} The greedy controller shows how large the density effect
        alone is; we did not run an experiment with an arena that grows with $N$.
  \item \textbf{``Architecture'' is a bundle.} The attention policy differs from the shared policy in set structure,
        attention weights, an explicit distance input and the number of parameters. We did not separate these factors,
        so we cannot say which one matters.
  \item \textbf{Reference runs.} The policies trained directly on four agents use the same 400 iterations as all others and
        one attention seed is clearly worse than the other two, so ``trained on 4'' is not a converged optimum.
  \item \textbf{Tuning.} The collision penalty was tuned on the same evaluation worlds with the shared policy
        (Section~\ref{sec:pt-training}); all other settings were fixed in advance.
  \item \textbf{Extension rule for independent policies.} Giving the new agent a copy of agent~1's network is one choice
        among several (random initialisation, an average of the trained networks); the results for independent
        policies depend on it.
  \item \textbf{Metrics.} Success is measured at the final step only, and evaluation uses greedy actions; stochastic
        execution was not studied.
\end{itemize}

\paragraph{Future work.}
(i) Train on a mixture of team sizes; (ii) isolate which components of the attention architecture matter;
(iii) test whether the findings persist with continuous dynamics, stochastic execution and larger arenas; and
(iv) combine the population axis with the spatial-scale axis studied in the other parts of the project.

\subsection{Reproducibility}
\label{sec:pt-repro}
All code is in the repository \texttt{marl-population-transfer} (files \texttt{env.py}, \texttt{policies.py},
\texttt{train.py}, \texttt{evaluate.py}, and \texttt{make\_figures.py}).
\texttt{run\_all.sh} trains all policies (three seeds, 400 iterations) and the four-agent references and produces
every table and figure of this section; it takes roughly 15--20 minutes on a laptop CPU. Evaluation uses a fixed
random seed (12345) so that all policies are tested on identical worlds, and every table entry in this section is
generated from the stored results file.

%  Needs these packages (ICLR style already loads most of them): amsmath, amssymb, amsthm,
%  booktabs, graphicx, natbib.  No tikz is needed: the architecture diagram is a PDF.
%  Figures (folder figures/): architectures.pdf, transfer_metrics.png, training_curves.png.
%  Citation keys are defined in refs_population_transfer.bib.
%  All numbers in the tables are generated from runs/results.json by the code repository.
% ==============================================================================================
\makeatletter
\@ifundefined{proposition}{\newtheorem{proposition}{Proposition}}{}
\@ifundefined{remark}{\theoremstyle{remark}\newtheorem{remark}{Remark}\theoremstyle{plain}}{}
\makeatother

\section{Population Transfer with Shared, Independent and Attention Policies}
\label{sec:pt}

\subsection{Overview and relation to the proposal}
\label{sec:pt-overview}
Our proposal asks whether a decentralised policy that satisfies a constraint on one configuration of agents
keeps satisfying it after the configuration changes. This section studies the \emph{population} axis of that
question, $N\to N{+}k$ agents, in the simplest setting in which it can be asked cleanly: cooperative
navigation with a collision-avoidance constraint. The world is deliberately small, so that every number can be
reproduced on a CPU in minutes and so that the effect of the \emph{policy architecture} on constraint
preservation is not hidden behind other design choices. 

\paragraph{Contributions.}
\begin{enumerate}
  \item A transfer protocol that judges a policy by its \emph{constraint-related} metrics and not by reward
        alone (Section~\ref{sec:pt-setting}).
  \item A comparison of three policy families that differ in how they handle a change in the number of agents:
        independent per-agent networks, a single shared network, and a shared attention network over neighbours,
        plus an ablation that separates architecture from information (Section~\ref{sec:pt-policies}).
  \item Empirical findings on how much of the constraint survives transfer and which claims our data do and do not
        support (Sections~\ref{sec:pt-results}--\ref{sec:pt-limits}).
\end{enumerate}

\paragraph{Metrics of the proposal.}
The proposal measures transfer through consensus error, minimum inter-agent separation and constraint-violation
rate. Our task is goal reaching and not consensus, so there is no consensus error here. We measure the other two
quantities, and we add goal-reaching success as the measure of task performance.

\subsection{Related work}
\label{sec:pt-related}
\paragraph{Independent learners and parameter sharing.}
Independent learners train one policy per agent \citep{tan1993}. Parameter sharing uses one network for all agents
and is the standard way to obtain a policy that is defined for any team size \citep{gupta2017}. Both are strong
baselines in cooperative tasks when trained with PPO \citep{dewitt2020,yu2022}. Centralised training with
decentralised execution \citep{lowe2017,foerster2016} is the setting assumed in our proposal; our agents are
trained with local critics and executed in a fully decentralised way.

\paragraph{Set-based and attention architectures.}
Mean-field methods \citep{yang2018}, attention-based critics \citep{iqbal2019} and permutation-invariant set
encoders \citep{zaheer2017} treat the neighbours of an agent as a set and therefore apply to any number of
neighbours. Our attention policy follows this idea with a single attention head \citep{vaswani2017}.

\paragraph{Transfer across agent populations.}
Transfer in reinforcement learning is surveyed by \citet{taylor2009}. Our focus is how a \emph{collision
constraint} survives a change in the \emph{number of agents}.

\subsection{Problem setting}
\label{sec:pt-setting}
\paragraph{World.}
$N$ agents move in the unit square $[0,1]^2$. Agent $i$ has position $p_i^t$ and a goal $g_i$. There are five
actions $\mathcal A=\{\text{stay},\text{up},\text{down},\text{left},\text{right}\}$ with displacements
$\delta(a)\in\{0,\pm s e_1,\pm s e_2\}$, so $\|\delta(a)\|\le s$ with step size $s=0.05$. Positions are updated
simultaneously,
\begin{equation}
  p_i^{t+1}=\Pi_{[0,1]^2}\!\big(p_i^t+\delta(a_i^t)\big),
\end{equation}
where $\Pi$ projects back into the arena. An episode lasts $T=50$ steps. Starting positions and goals are
sampled uniformly, rejecting samples in which two starts or two goals are closer than $0.15$.

\paragraph{Constraint.}
Two agents \emph{collide} when $\|p_i^t-p_j^t\|<d_c$ with $d_c=0.1$. The constraint is
\begin{equation}
  \|p_i^t-p_j^t\|\ \ge\ d_c\qquad\text{for all } i\neq j \text{ and all } t.
\end{equation}

\paragraph{Observations.}
Every agent acts only on local information. With $r_{ij}=p_j-p_i$ and $\mathcal N_i^k$ the $k$ agents nearest to $i$,
\begin{equation}
  o_i^t=\big(p_i^t,\; g_i-p_i^t,\; (r_{ij}^t)_{j\in\mathcal N_i^k}\big)\in\mathbb R^{4+2k}.
\end{equation}
The shared and independent policies use $k=2$, so the observation length does not depend on $N$; this is what
allows a network trained with $N$ agents to be executed with $N{+}1$. The attention policy receives the
neighbours as a set whose size may vary (all $N{-}1$ other agents, or only the two nearest in the ablation).

\paragraph{Reward.}
Each agent receives
\begin{equation}
  r_i^t=10\big(\|g_i-p_i^{t-1}\|-\|g_i-p_i^t\|\big)\;-\;\lambda\,\mathbf{1}\!\big[\min_{j\neq i}\|p_i^t-p_j^t\|<d_c\big]\;+\;0.2\,\mathbf{1}\!\big[\|g_i-p_i^t\|<0.05\big],
\end{equation}
with collision penalty $\lambda=10$. The collision term is the only coupling between agents.

\paragraph{Transfer protocol.}
Each policy is trained on $N_{\text{train}}=3$ agents. It is then evaluated \emph{without any retraining} on
$N\in\{3,4,5,6\}$ agents in the same arena, so the team becomes denser as $N$ grows. As a reference we also train
each policy family directly on 4 agents. The difference between the transferred and the directly trained policy
at 4 agents is the \emph{transfer gap} (Section~\ref{sec:pt-results}).

\paragraph{Metrics.}
For each setting we run 500 evaluation episodes on the same random worlds for all policies (a paired comparison)
and use greedy (arg\,max) actions. We report means and standard deviations over three training seeds:
\begin{itemize}
  \item \textbf{Minimum separation:} $\min_{t,\,i\neq j}\|p_i^t-p_j^t\|$ within an episode, averaged over episodes
        (higher is safer; values below $d_c$ mean at least one collision).
  \item \textbf{Violation rate:} $\frac1T\sum_t\mathbf{1}[\exists\, i\neq j:\ \|p_i^t-p_j^t\|<d_c]$, the fraction of
        time steps with a collision.
  \item \textbf{Episodes with a collision:} the fraction of episodes with at least one collision.
  \item \textbf{Success rate:} the fraction of agents with $\|g_i-p_i^T\|<0.05$ at the final step
        (task performance).
\end{itemize}
A hand-coded \emph{greedy} controller that walks straight to its goal and ignores all other agents is included as a
``no avoidance'' reference. It shows how much of the loss of safety is due to the increased density alone.

\subsection{Policy architectures}
\label{sec:pt-policies}
\begin{figure}[t]
  \centering
  \includegraphics[width=0.78\linewidth]{figures/architectures.pdf}
  \caption{The three policy families, shown for a team trained on three agents that gets a fourth agent at test
  time (dashed). (a)~One network is shared by all agents. (b)~Each agent has its own network; the new agent has
  none and borrows a copy of agent~1's network. (c)~One shared network encodes the neighbours of an agent as a set
  and combines them with attention, so an extra neighbour only adds one element to the set.}
  \label{fig:pt-arch}
\end{figure}

\paragraph{Shared policy.}
One actor $\pi_\theta(a\mid o_i)$ and one critic $V_\phi(o_i)$, each a multilayer perceptron with two hidden layers
of 64 $\tanh$ units, are applied to every agent \citep{gupta2017}. An additional agent simply uses the same
network, so the policy is defined for every $N$ by construction.

\paragraph{Independent policies.}
Agent $i$ has its own actor $\pi_{\theta_i}$ and critic $V_{\phi_i}$ \citep{tan1993}. For $N>N_{\text{train}}$
there is no trained network for the additional agents, so an extension rule is needed. We let agent $i$ use the
network $\sigma(i)=((i-1)\bmod N_{\text{train}})+1$, i.e.\ the fourth agent receives a copy of agent~1's network.
This is a design choice and not a property of the method (Section~\ref{sec:pt-limits}).

\paragraph{Attention policy.}
With $r_{ij}=p_j-p_i$, the policy of agent $i$ computes
\begin{align}
  h_i &= \tanh\!\big(W_s[\,p_i;\ g_i-p_i\,]+b_s\big), &
  e_{ij} &= \tanh\!\big(W_n[\,r_{ij};\ \|r_{ij}\|\,]+b_n\big),\\
  \alpha_{ij} &= \operatorname*{softmax}_{j}\ \frac{(W_q h_i)^{\top}(W_k e_{ij})}{\sqrt d}, &
  c_i &= \sum_{j\neq i}\alpha_{ij}\,W_v e_{ij},
\end{align}
and outputs $\pi(\cdot\mid o_i)=\operatorname{softmax}\big(f_\theta([h_i;c_i])\big)$ with $d=64$ and $f_\theta$ a
two-layer MLP; the critic has the same structure. Because $c_i$ is a weighted average over the neighbours, it has
a fixed size and does not depend on the order of the neighbours: the network is permutation invariant and accepts
any number of neighbours \citep{zaheer2017}.

\paragraph{Ablation: attention on the two nearest neighbours.}
The attention policy sees all other agents, whereas the shared and independent policies see only the two nearest.
To separate architecture from information we also train the \emph{same} attention network on the two nearest
neighbours only. With three agents the two variants receive identical observations, so they coincide at $N=3$ by
construction; they can differ only for $N>3$.

\begin{table}[t]
  \centering\small
  \caption{Summary of the policies. ``Parameters'' counts actor and critic together for a team trained on three agents.}
  \label{tab:pt-arch}
  \begin{tabular}{l c l p{4.1cm} r}
    \toprule
    Policy & Networks & Neighbours seen & Extra agent at test time & Parameters \\
    \midrule
    Independent            & $N_{\text{train}}$ & 2 nearest & copy of the network of agent $\sigma(i)$ & 29,586 \\
    Shared                 & 1 & 2 nearest & same network & 9,862 \\
    Attention (all agents) & 1 & all other agents & same network & 38,278 \\
    Attention (2 nearest)  & 1 & 2 nearest & same network & 38,278 \\
    \bottomrule
  \end{tabular}
\end{table}

\subsection{Training}
\label{sec:pt-training}
All policies are trained with proximal policy optimisation \citep{schulman2017} and generalised advantage
estimation \citep{schulman2016}. Every agent uses only its own observation, both for the actor and for a local
critic, which corresponds to independent PPO \citep{dewitt2020}; only the actor is needed at execution time. One
iteration collects 50-step episodes in 64 parallel worlds and then updates the networks for four epochs. For the
independent policies the loss of every agent's network is computed on that agent's own experience. The settings are
listed in Table~\ref{tab:pt-hyper}.

\begin{table}[t]
  \centering\small
  \caption{Training settings (identical for all policy families).}
  \label{tab:pt-hyper}
  \begin{tabular}{ll}
    \toprule
    Quantity & Value \\
    \midrule
    Algorithm & PPO with GAE, local critics \\
    Optimiser & Adam, learning rate $10^{-3}$ \\
    Discount $\gamma$ / GAE $\lambda_{\text{GAE}}$ & 0.97 / 0.95 \\
    Clip range / entropy coefficient / value coefficient & 0.2 / 0.01 / 0.5 \\
    Gradient-norm clip & 0.5 \\
    Epochs per iteration / minibatch (world steps) & 4 / 512 \\
    Parallel worlds $\times$ steps per iteration & $64\times 50$ \\
    Iterations & 400 ($\approx$ 1.3 million world steps) \\
    Reward weights (progress / collision / goal) & 10 / 10 / 0.2 \\
    Training seeds & 0, 1, 2 \\
    \bottomrule
  \end{tabular}
\end{table}

\paragraph{The one tuned quantity.}
The collision penalty $\lambda$ was the only value we tuned. In a pilot with the shared policy, $\lambda=1$ and
$\lambda=3$ produced behaviour that was close to the greedy controller (agents had little incentive to avoid each
other), and $\lambda=10$ was the next value we tried and the first to give clear avoidance. We then used
$\lambda=10$ for all policies and all experiments. The pilot used the same evaluation worlds as the final
evaluation, which can make the absolute collision numbers slightly optimistic; it cannot favour one policy family,
since the value was fixed before the other families were trained.

\subsection{Experiments and results}
\label{sec:pt-results}
We train three seeds of every policy family on three agents (and, for the reference, three seeds of each of the
three main families on four agents) and evaluate on 3, 4, 5 and 6 agents.
All tables give the mean $\pm$ standard deviation over the three training seeds (population standard deviation).
With only three seeds we interpret large differences and treat differences of the order of the seed spread as
inconclusive.

\begin{figure}[t]
  \centering
  \includegraphics[width=\linewidth]{figures/transfer_metrics.png}
  \caption{Transfer from three to $N$ agents without retraining. Lines are means over three training seeds, faint dots
  are individual seeds, and the hollow diamonds are policies trained directly on four agents. \emph{Left:} minimum
  separation (collision below $0.1$). \emph{Middle:} fraction of time steps with a collision. \emph{Right:} fraction
  of agents on their goal at the end. The ``2 nearest'' attention variant (triangles) lies on top of the attention
  policy that sees all agents.}
  \label{fig:pt-transfer}
\end{figure}

\begin{table}[t]
  \centering\small
  \caption{Percentage of episodes with at least one collision (lower is safer).}
  \label{tab:pt-collision}
\begin{tabular}{l cccc}
\toprule
 & \multicolumn{4}{c}{Number of agents at test time} \\
\cmidrule(lr){2-5}
Policy (trained on 3 agents) & 3 (trained) & 4 & 5 & 6 \\
\midrule
Greedy (no avoidance) & 40.0 $\pm$ 0.0 & 66.6 $\pm$ 0.0 & 81.6 $\pm$ 0.0 & 92.8 $\pm$ 0.0 \\
\midrule
Independent & 20.7 $\pm$ 1.9 & 38.7 $\pm$ 1.4 & 58.4 $\pm$ 2.6 & 71.7 $\pm$ 2.2 \\
Shared & 16.3 $\pm$ 2.0 & 32.8 $\pm$ 3.9 & 49.5 $\pm$ 4.2 & 66.1 $\pm$ 3.4 \\
Attention (all agents) & 1.0 $\pm$ 0.2 & 3.0 $\pm$ 0.8 & 9.0 $\pm$ 3.1 & 15.3 $\pm$ 4.2 \\
Attention (2 nearest) & 1.0 $\pm$ 0.2 & 2.9 $\pm$ 1.1 & 8.3 $\pm$ 3.6 & 12.8 $\pm$ 5.7 \\
\midrule
\bottomrule
\end{tabular}
\end{table}

\begin{table}[t]
  \centering\small
  \caption{Minimum separation within an episode, averaged over episodes (higher is safer; the collision distance is
  $0.1$).}
  \label{tab:pt-minsep}
\begin{tabular}{l cccc}
\toprule
 & \multicolumn{4}{c}{Number of agents at test time} \\
\cmidrule(lr){2-5}
Policy (trained on 3 agents) & 3 (trained) & 4 & 5 & 6 \\
\midrule
Greedy (no avoidance) & 0.144 $\pm$ 0.000 & 0.089 $\pm$ 0.000 & 0.061 $\pm$ 0.000 & 0.045 $\pm$ 0.000 \\
\midrule
Independent & 0.192 $\pm$ 0.005 & 0.132 $\pm$ 0.003 & 0.095 $\pm$ 0.003 & 0.073 $\pm$ 0.004 \\
Shared & 0.201 $\pm$ 0.005 & 0.139 $\pm$ 0.005 & 0.104 $\pm$ 0.005 & 0.082 $\pm$ 0.004 \\
Attention (all agents) & 0.214 $\pm$ 0.008 & 0.172 $\pm$ 0.008 & 0.148 $\pm$ 0.006 & 0.131 $\pm$ 0.004 \\
Attention (2 nearest) & 0.214 $\pm$ 0.008 & 0.173 $\pm$ 0.008 & 0.150 $\pm$ 0.007 & 0.134 $\pm$ 0.007 \\
\midrule
\bottomrule
\end{tabular}
\end{table}

\begin{table}[t]
  \centering\small
  \caption{Percentage of time steps with a collision.}
  \label{tab:pt-violation}
\begin{tabular}{l cccc}
\toprule
 & \multicolumn{4}{c}{Number of agents at test time} \\
\cmidrule(lr){2-5}
Policy (trained on 3 agents) & 3 (trained) & 4 & 5 & 6 \\
\midrule
Greedy (no avoidance) & 2.80 $\pm$ 0.00 & 5.61 $\pm$ 0.00 & 8.76 $\pm$ 0.00 & 12.74 $\pm$ 0.00 \\
\midrule
Independent & 1.01 $\pm$ 0.11 & 2.07 $\pm$ 0.11 & 3.83 $\pm$ 0.25 & 5.64 $\pm$ 0.48 \\
Shared & 0.76 $\pm$ 0.11 & 1.71 $\pm$ 0.14 & 3.23 $\pm$ 0.09 & 4.96 $\pm$ 0.15 \\
Attention (all agents) & 0.04 $\pm$ 0.02 & 0.17 $\pm$ 0.09 & 0.42 $\pm$ 0.19 & 0.90 $\pm$ 0.31 \\
Attention (2 nearest) & 0.04 $\pm$ 0.02 & 0.16 $\pm$ 0.11 & 0.42 $\pm$ 0.23 & 0.95 $\pm$ 0.52 \\
\bottomrule
\end{tabular}
\end{table}

\begin{table}[t]
  \centering\small
  \caption{Success rate in percent (agents on their goal at the end of the episode).}
  \label{tab:pt-success}
\begin{tabular}{l cccc}
\toprule
 & \multicolumn{4}{c}{Number of agents at test time} \\
\cmidrule(lr){2-5}
Policy (trained on 3 agents) & 3 (trained) & 4 & 5 & 6 \\
\midrule
Greedy (no avoidance) & 100.0 $\pm$ 0.0 & 100.0 $\pm$ 0.0 & 100.0 $\pm$ 0.0 & 100.0 $\pm$ 0.0 \\
\midrule
Independent & 96.2 $\pm$ 0.6 & 96.3 $\pm$ 0.3 & 97.0 $\pm$ 0.2 & 97.1 $\pm$ 0.7 \\
Shared & 98.8 $\pm$ 0.5 & 98.9 $\pm$ 0.5 & 98.9 $\pm$ 0.6 & 98.8 $\pm$ 0.7 \\
Attention (all agents) & 97.4 $\pm$ 1.6 & 96.5 $\pm$ 1.6 & 94.9 $\pm$ 1.4 & 94.1 $\pm$ 1.3 \\
Attention (2 nearest) & 97.4 $\pm$ 1.6 & 96.2 $\pm$ 2.0 & 94.5 $\pm$ 1.9 & 92.9 $\pm$ 2.2 \\
\midrule
\bottomrule
\end{tabular}
\end{table}

\paragraph{Task success transfers, constraint satisfaction does not.}
Every learned policy reaches its goals in at least 92.9\,\% of the cases at every team size
(Table~\ref{tab:pt-success}), so the task itself transfers. The constraint does not. The share of episodes with a
collision (Table~\ref{tab:pt-collision}) is $16.3\,\%$ (shared) and $20.7\,\%$ (independent) with three agents and
roughly doubles when a fourth agent is added ($32.8\,\%$ and $38.7\,\%$); with six agents it is $66.1\,\%$ and
$71.7\,\%$. The mean minimum separation of the independent policies falls below the collision distance of $0.1$ at five
agents and that of the shared policy at six agents (Table~\ref{tab:pt-minsep}). A policy that looks successful by
its task reward can therefore already have lost much of its safety.

\paragraph{Part of the loss is density.}
The greedy controller, which has no avoidance at all, goes from $40.0\,\%$ to $92.8\,\%$ of episodes with a collision
between three and six agents, so a denser team is harder for any controller. The shared and independent policies
stay below the greedy controller at every size, but their advantage over it stays between $19$ and $34$ percentage points
and is smaller with six agents than with four: the increase from three to six agents is $49.8$ and $50.9$ points for
the shared and independent policies, against $52.8$ for the greedy controller. The advantage of the attention policy
over the greedy controller, in contrast, keeps growing, from $39$ to $78$ points.

\paragraph{The attention policy degrades much more slowly, and the cause is the architecture.}
The attention policy has $1.0\,\%$ of episodes with a collision at three agents and $15.3\,\%$ at six (an increase of
$14.3$ points), with a mean minimum separation of $0.131$ at six agents, still above $d_c$.
One might suspect that this is only because it sees all agents, whereas the other two see the two nearest. The
ablation rules this out for our setting: the same network fed only the two nearest neighbours reaches $12.8\,\%$ at
six agents (versus $15.3\,\%$; the difference is smaller than the seed spread of $4$--$6$ points) and is
indistinguishable at four and five agents (Table~\ref{tab:pt-collision}, Figure~\ref{fig:pt-transfer}). Seeing more
neighbours gave no measurable benefit. The comparison between the attention-on-two-nearest policy and the shared
policy differs in the architecture only, but ``architecture'' here is a bundle: set-structure with attention weights, an
explicit distance input $\|r_{ij}\|$, separate encoders, and about four times more parameters (Table~\ref{tab:pt-arch}).
We did not ablate these factors individually.

\paragraph{Shared versus independent.}
The shared policy is better than the independent policies in every column: $4.4$ to $8.9$ points fewer episodes with a
collision, a larger minimum separation and a higher success rate ($98.8\,\%$ versus $96.2\,\%$ at three agents). The
collision differences are of the same order as the seed spread ($1.4$--$4.2$ points), so we do not claim them to be
significant; the success difference is larger than the spread. Training curves
(Figure~\ref{fig:pt-train}) show that the shared policy learns faster than the independent ones, which is expected
since one network learns from the experience of all agents, and that the attention policy learns fastest.

\paragraph{Transfer gap.}
Table~\ref{tab:pt-gap} compares policies transferred from three to four agents with policies trained directly on four
agents. For the shared and independent policies the gap is large: $+17.7$ and $+12.9$ percentage points of episodes with a
collision. Directly trained policies are safer but succeed slightly less often, i.e.\ they trade some task success for
safety. For the attention policy there is no measurable gap ($-5.3$ points with a standard deviation of $8.7$); the
directly trained attention policies have $20.6\,\%$, $2.6\,\%$ and $1.8\,\%$ in the three seeds, so the poor average
comes from a single seed, and the transferred policy ($3.0\,\%$) is as safe as the two good ones. The reference runs use
the same training budget as the others and are therefore not necessarily converged.

\begin{table}[t]
  \centering\small
  \caption{Transfer gap at four agents: policies trained on three agents and transferred, versus policies trained
  directly on four agents. A positive gap means the transferred policy has more collisions (percentage points).}
  \label{tab:pt-gap}
\begin{tabular}{l cc c cc}
\toprule
 & \multicolumn{3}{c}{Episodes with a collision (\%)} & \multicolumn{2}{c}{Success rate (\%)} \\
\cmidrule(lr){2-4}\cmidrule(lr){5-6}
Policy & Transferred ($3\!\to\!4$) & Trained on 4 & Gap (pp) & Transferred & Trained on 4 \\
\midrule
Independent & 38.7 $\pm$ 1.4 & 25.8 $\pm$ 1.1 & +12.9 & 96.3 $\pm$ 0.3 & 92.0 $\pm$ 0.5 \\
Shared & 32.8 $\pm$ 3.9 & 15.1 $\pm$ 6.3 & +17.7 & 98.9 $\pm$ 0.5 & 96.9 $\pm$ 1.3 \\
Attention (all agents) & 3.0 $\pm$ 0.8 & 8.3 $\pm$ 8.7 & -5.3 & 96.5 $\pm$ 1.6 & 93.7 $\pm$ 4.0 \\
\bottomrule
\end{tabular}
\end{table}

\begin{figure}[t]
  \centering
  \includegraphics[width=0.92\linewidth]{figures/training_curves.png}
  \caption{Training curves on three agents (mean over three seeds, band = min to max; moving average of 10
  iterations). Left: training return. Right: fraction of time steps with a collision during training.}
  \label{fig:pt-train}
\end{figure}

\subsection{Discussion}
\label{sec:pt-discussion}
\paragraph{What this says about constraint-preserving transfer.}
Transfer of the task and transfer of the constraint are different things. The task survived a change of population
almost perfectly, while the constraint did not, and the loss was not predicted by task performance. For the
proposal this means that preservation of a constraint cannot be assumed from successful transfer; it has to be
measured directly. In our experiments, the set-based attention architecture preserved much more of the constraint
than the concatenated-observation policies, and its transfer gap at four agents was not measurable.

\paragraph{Information.}
When the team outgrows a $k$-nearest observation, each agent no longer observes every other agent. Our ablation
shows that, for the attention policy, observing all neighbours rather than only the two nearest did not improve the
measured collision outcomes. In this setting, the transfer gap is therefore not explained by this information
difference alone.

\paragraph{Practical reading.}
If a team may grow, the set-based attention policy is the most robust of the learned policies we tested: it has the
lowest collision rate while retaining high goal-reaching success across all evaluated team sizes.

\subsection{Limitations and threats to validity}
\label{sec:pt-limits}
\begin{itemize}
  \item \textbf{A small world.} The dynamics are discrete and noise free, there is no communication, and the arena is
        fixed, so the team becomes denser as it grows. Conclusions about larger teams, continuous dynamics or noise are
        not supported.
  \item \textbf{Three seeds.} We report means and standard deviations over three seeds without confidence intervals;
        only large effects (for example attention versus the other policy families) are reliable. The comparison between
        shared and independent policies is suggestive only.
  \item \textbf{Density and policy are confounded in $N$.} The greedy controller shows how large the density effect
        alone is; we did not run an experiment with an arena that grows with $N$.
  \item \textbf{``Architecture'' is a bundle.} The attention policy differs from the shared policy in set structure,
        attention weights, an explicit distance input and the number of parameters. We did not separate these factors,
        so we cannot say which one matters.
  \item \textbf{Reference runs.} The policies trained directly on four agents use the same 400 iterations as all others and
        one attention seed is clearly worse than the other two, so ``trained on 4'' is not a converged optimum.
  \item \textbf{Tuning.} The collision penalty was tuned on the same evaluation worlds with the shared policy
        (Section~\ref{sec:pt-training}); all other settings were fixed in advance.
  \item \textbf{Extension rule for independent policies.} Giving the new agent a copy of agent~1's network is one choice
        among several (random initialisation, an average of the trained networks); the results for independent
        policies depend on it.
  \item \textbf{Metrics.} Success is measured at the final step only, and evaluation uses greedy actions; stochastic
        execution was not studied.
\end{itemize}

\paragraph{Future work.}
(i) Train on a mixture of team sizes; (ii) isolate which components of the attention architecture matter;
(iii) test whether the findings persist with continuous dynamics, stochastic execution and larger arenas; and
(iv) combine the population axis with the spatial-scale axis studied in the other parts of the project.

\subsection{Reproducibility}
\label{sec:pt-repro}
All code is in the repository \texttt{marl-population-transfer} (files \texttt{env.py}, \texttt{policies.py},
\texttt{train.py}, \texttt{evaluate.py}, and \texttt{make\_figures.py}).
\texttt{run\_all.sh} trains all policies (three seeds, 400 iterations) and the four-agent references and produces
every table and figure of this section; it takes roughly 15--20 minutes on a laptop CPU. Evaluation uses a fixed
random seed (12345) so that all policies are tested on identical worlds, and every table entry in this section is
generated from the stored results file.

%  Needs these packages (ICLR style already loads most of them): amsmath, amssymb, amsthm,
%  booktabs, graphicx, natbib.  No tikz is needed: the architecture diagram is a PDF.
%  Figures (folder figures/): architectures.pdf, transfer_metrics.png, training_curves.png.
%  Citation keys are defined in refs_population_transfer.bib.
%  All numbers in the tables are generated from runs/results.json by the code repository.
% ==============================================================================================
\makeatletter
\@ifundefined{proposition}{\newtheorem{proposition}{Proposition}}{}
\@ifundefined{remark}{\theoremstyle{remark}\newtheorem{remark}{Remark}\theoremstyle{plain}}{}
\makeatother

\section{Population Transfer with Shared, Independent and Attention Policies}
\label{sec:pt}

\subsection{Overview and relation to the proposal}
\label{sec:pt-overview}
Our proposal asks whether a decentralised policy that satisfies a constraint on one configuration of agents
keeps satisfying it after the configuration changes. This section studies the \emph{population} axis of that
question, $N\to N{+}k$ agents, in the simplest setting in which it can be asked cleanly: cooperative
navigation with a collision-avoidance constraint. The world is deliberately small, so that every number can be
reproduced on a CPU in minutes and so that the effect of the \emph{policy architecture} on constraint
preservation is not hidden behind other design choices. 

\paragraph{Contributions.}
\begin{enumerate}
  \item A transfer protocol that judges a policy by its \emph{constraint-related} metrics and not by reward
        alone (Section~\ref{sec:pt-setting}).
  \item A comparison of three policy families that differ in how they handle a change in the number of agents:
        independent per-agent networks, a single shared network, and a shared attention network over neighbours,
        plus an ablation that separates architecture from information (Section~\ref{sec:pt-policies}).
  \item Empirical findings on how much of the constraint survives transfer and which claims our data do and do not
        support (Sections~\ref{sec:pt-results}--\ref{sec:pt-limits}).
\end{enumerate}

\paragraph{Metrics of the proposal.}
The proposal measures transfer through consensus error, minimum inter-agent separation and constraint-violation
rate. Our task is goal reaching and not consensus, so there is no consensus error here. We measure the other two
quantities, and we add goal-reaching success as the measure of task performance.

\subsection{Related work}
\label{sec:pt-related}
\paragraph{Independent learners and parameter sharing.}
Independent learners train one policy per agent \citep{tan1993}. Parameter sharing uses one network for all agents
and is the standard way to obtain a policy that is defined for any team size \citep{gupta2017}. Both are strong
baselines in cooperative tasks when trained with PPO \citep{dewitt2020,yu2022}. Centralised training with
decentralised execution \citep{lowe2017,foerster2016} is the setting assumed in our proposal; our agents are
trained with local critics and executed in a fully decentralised way.

\paragraph{Set-based and attention architectures.}
Mean-field methods \citep{yang2018}, attention-based critics \citep{iqbal2019} and permutation-invariant set
encoders \citep{zaheer2017} treat the neighbours of an agent as a set and therefore apply to any number of
neighbours. Our attention policy follows this idea with a single attention head \citep{vaswani2017}.

\paragraph{Transfer across agent populations.}
Transfer in reinforcement learning is surveyed by \citet{taylor2009}. Our focus is how a \emph{collision
constraint} survives a change in the \emph{number of agents}.

\subsection{Problem setting}
\label{sec:pt-setting}
\paragraph{World.}
$N$ agents move in the unit square $[0,1]^2$. Agent $i$ has position $p_i^t$ and a goal $g_i$. There are five
actions $\mathcal A=\{\text{stay},\text{up},\text{down},\text{left},\text{right}\}$ with displacements
$\delta(a)\in\{0,\pm s e_1,\pm s e_2\}$, so $\|\delta(a)\|\le s$ with step size $s=0.05$. Positions are updated
simultaneously,
\begin{equation}
  p_i^{t+1}=\Pi_{[0,1]^2}\!\big(p_i^t+\delta(a_i^t)\big),
\end{equation}
where $\Pi$ projects back into the arena. An episode lasts $T=50$ steps. Starting positions and goals are
sampled uniformly, rejecting samples in which two starts or two goals are closer than $0.15$.

\paragraph{Constraint.}
Two agents \emph{collide} when $\|p_i^t-p_j^t\|<d_c$ with $d_c=0.1$. The constraint is
\begin{equation}
  \|p_i^t-p_j^t\|\ \ge\ d_c\qquad\text{for all } i\neq j \text{ and all } t.
\end{equation}

\paragraph{Observations.}
Every agent acts only on local information. With $r_{ij}=p_j-p_i$ and $\mathcal N_i^k$ the $k$ agents nearest to $i$,
\begin{equation}
  o_i^t=\big(p_i^t,\; g_i-p_i^t,\; (r_{ij}^t)_{j\in\mathcal N_i^k}\big)\in\mathbb R^{4+2k}.
\end{equation}
The shared and independent policies use $k=2$, so the observation length does not depend on $N$; this is what
allows a network trained with $N$ agents to be executed with $N{+}1$. The attention policy receives the
neighbours as a set whose size may vary (all $N{-}1$ other agents, or only the two nearest in the ablation).

\paragraph{Reward.}
Each agent receives
\begin{equation}
  r_i^t=10\big(\|g_i-p_i^{t-1}\|-\|g_i-p_i^t\|\big)\;-\;\lambda\,\mathbf{1}\!\big[\min_{j\neq i}\|p_i^t-p_j^t\|<d_c\big]\;+\;0.2\,\mathbf{1}\!\big[\|g_i-p_i^t\|<0.05\big],
\end{equation}
with collision penalty $\lambda=10$. The collision term is the only coupling between agents.

\paragraph{Transfer protocol.}
Each policy is trained on $N_{\text{train}}=3$ agents. It is then evaluated \emph{without any retraining} on
$N\in\{3,4,5,6\}$ agents in the same arena, so the team becomes denser as $N$ grows. As a reference we also train
each policy family directly on 4 agents. The difference between the transferred and the directly trained policy
at 4 agents is the \emph{transfer gap} (Section~\ref{sec:pt-results}).

\paragraph{Metrics.}
For each setting we run 500 evaluation episodes on the same random worlds for all policies (a paired comparison)
and use greedy (arg\,max) actions. We report means and standard deviations over three training seeds:
\begin{itemize}
  \item \textbf{Minimum separation:} $\min_{t,\,i\neq j}\|p_i^t-p_j^t\|$ within an episode, averaged over episodes
        (higher is safer; values below $d_c$ mean at least one collision).
  \item \textbf{Violation rate:} $\frac1T\sum_t\mathbf{1}[\exists\, i\neq j:\ \|p_i^t-p_j^t\|<d_c]$, the fraction of
        time steps with a collision.
  \item \textbf{Episodes with a collision:} the fraction of episodes with at least one collision.
  \item \textbf{Success rate:} the fraction of agents with $\|g_i-p_i^T\|<0.05$ at the final step
        (task performance).
\end{itemize}
A hand-coded \emph{greedy} controller that walks straight to its goal and ignores all other agents is included as a
``no avoidance'' reference. It shows how much of the loss of safety is due to the increased density alone.

\subsection{Policy architectures}
\label{sec:pt-policies}
\begin{figure}[t]
  \centering
  \includegraphics[width=0.78\linewidth]{figures/architectures.pdf}
  \caption{The three policy families, shown for a team trained on three agents that gets a fourth agent at test
  time (dashed). (a)~One network is shared by all agents. (b)~Each agent has its own network; the new agent has
  none and borrows a copy of agent~1's network. (c)~One shared network encodes the neighbours of an agent as a set
  and combines them with attention, so an extra neighbour only adds one element to the set.}
  \label{fig:pt-arch}
\end{figure}

\paragraph{Shared policy.}
One actor $\pi_\theta(a\mid o_i)$ and one critic $V_\phi(o_i)$, each a multilayer perceptron with two hidden layers
of 64 $\tanh$ units, are applied to every agent \citep{gupta2017}. An additional agent simply uses the same
network, so the policy is defined for every $N$ by construction.

\paragraph{Independent policies.}
Agent $i$ has its own actor $\pi_{\theta_i}$ and critic $V_{\phi_i}$ \citep{tan1993}. For $N>N_{\text{train}}$
there is no trained network for the additional agents, so an extension rule is needed. We let agent $i$ use the
network $\sigma(i)=((i-1)\bmod N_{\text{train}})+1$, i.e.\ the fourth agent receives a copy of agent~1's network.
This is a design choice and not a property of the method (Section~\ref{sec:pt-limits}).

\paragraph{Attention policy.}
With $r_{ij}=p_j-p_i$, the policy of agent $i$ computes
\begin{align}
  h_i &= \tanh\!\big(W_s[\,p_i;\ g_i-p_i\,]+b_s\big), &
  e_{ij} &= \tanh\!\big(W_n[\,r_{ij};\ \|r_{ij}\|\,]+b_n\big),\\
  \alpha_{ij} &= \operatorname*{softmax}_{j}\ \frac{(W_q h_i)^{\top}(W_k e_{ij})}{\sqrt d}, &
  c_i &= \sum_{j\neq i}\alpha_{ij}\,W_v e_{ij},
\end{align}
and outputs $\pi(\cdot\mid o_i)=\operatorname{softmax}\big(f_\theta([h_i;c_i])\big)$ with $d=64$ and $f_\theta$ a
two-layer MLP; the critic has the same structure. Because $c_i$ is a weighted average over the neighbours, it has
a fixed size and does not depend on the order of the neighbours: the network is permutation invariant and accepts
any number of neighbours \citep{zaheer2017}.

\paragraph{Ablation: attention on the two nearest neighbours.}
The attention policy sees all other agents, whereas the shared and independent policies see only the two nearest.
To separate architecture from information we also train the \emph{same} attention network on the two nearest
neighbours only. With three agents the two variants receive identical observations, so they coincide at $N=3$ by
construction; they can differ only for $N>3$.

\begin{table}[t]
  \centering\small
  \caption{Summary of the policies. ``Parameters'' counts actor and critic together for a team trained on three agents.}
  \label{tab:pt-arch}
  \begin{tabular}{l c l p{4.1cm} r}
    \toprule
    Policy & Networks & Neighbours seen & Extra agent at test time & Parameters \\
    \midrule
    Independent            & $N_{\text{train}}$ & 2 nearest & copy of the network of agent $\sigma(i)$ & 29,586 \\
    Shared                 & 1 & 2 nearest & same network & 9,862 \\
    Attention (all agents) & 1 & all other agents & same network & 38,278 \\
    Attention (2 nearest)  & 1 & 2 nearest & same network & 38,278 \\
    \bottomrule
  \end{tabular}
\end{table}

\subsection{Training}
\label{sec:pt-training}
All policies are trained with proximal policy optimisation \citep{schulman2017} and generalised advantage
estimation \citep{schulman2016}. Every agent uses only its own observation, both for the actor and for a local
critic, which corresponds to independent PPO \citep{dewitt2020}; only the actor is needed at execution time. One
iteration collects 50-step episodes in 64 parallel worlds and then updates the networks for four epochs. For the
independent policies the loss of every agent's network is computed on that agent's own experience. The settings are
listed in Table~\ref{tab:pt-hyper}.

\begin{table}[t]
  \centering\small
  \caption{Training settings (identical for all policy families).}
  \label{tab:pt-hyper}
  \begin{tabular}{ll}
    \toprule
    Quantity & Value \\
    \midrule
    Algorithm & PPO with GAE, local critics \\
    Optimiser & Adam, learning rate $10^{-3}$ \\
    Discount $\gamma$ / GAE $\lambda_{\text{GAE}}$ & 0.97 / 0.95 \\
    Clip range / entropy coefficient / value coefficient & 0.2 / 0.01 / 0.5 \\
    Gradient-norm clip & 0.5 \\
    Epochs per iteration / minibatch (world steps) & 4 / 512 \\
    Parallel worlds $\times$ steps per iteration & $64\times 50$ \\
    Iterations & 400 ($\approx$ 1.3 million world steps) \\
    Reward weights (progress / collision / goal) & 10 / 10 / 0.2 \\
    Training seeds & 0, 1, 2 \\
    \bottomrule
  \end{tabular}
\end{table}

\paragraph{The one tuned quantity.}
The collision penalty $\lambda$ was the only value we tuned. In a pilot with the shared policy, $\lambda=1$ and
$\lambda=3$ produced behaviour that was close to the greedy controller (agents had little incentive to avoid each
other), and $\lambda=10$ was the next value we tried and the first to give clear avoidance. We then used
$\lambda=10$ for all policies and all experiments. The pilot used the same evaluation worlds as the final
evaluation, which can make the absolute collision numbers slightly optimistic; it cannot favour one policy family,
since the value was fixed before the other families were trained.

\subsection{Experiments and results}
\label{sec:pt-results}
We train three seeds of every policy family on three agents (and, for the reference, three seeds of each of the
three main families on four agents) and evaluate on 3, 4, 5 and 6 agents.
All tables give the mean $\pm$ standard deviation over the three training seeds (population standard deviation).
With only three seeds we interpret large differences and treat differences of the order of the seed spread as
inconclusive.

\begin{figure}[t]
  \centering
  \includegraphics[width=\linewidth]{figures/transfer_metrics.png}
  \caption{Transfer from three to $N$ agents without retraining. Lines are means over three training seeds, faint dots
  are individual seeds, and the hollow diamonds are policies trained directly on four agents. \emph{Left:} minimum
  separation (collision below $0.1$). \emph{Middle:} fraction of time steps with a collision. \emph{Right:} fraction
  of agents on their goal at the end. The ``2 nearest'' attention variant (triangles) lies on top of the attention
  policy that sees all agents.}
  \label{fig:pt-transfer}
\end{figure}

\begin{table}[t]
  \centering\small
  \caption{Percentage of episodes with at least one collision (lower is safer).}
  \label{tab:pt-collision}
\begin{tabular}{l cccc}
\toprule
 & \multicolumn{4}{c}{Number of agents at test time} \\
\cmidrule(lr){2-5}
Policy (trained on 3 agents) & 3 (trained) & 4 & 5 & 6 \\
\midrule
Greedy (no avoidance) & 40.0 $\pm$ 0.0 & 66.6 $\pm$ 0.0 & 81.6 $\pm$ 0.0 & 92.8 $\pm$ 0.0 \\
\midrule
Independent & 20.7 $\pm$ 1.9 & 38.7 $\pm$ 1.4 & 58.4 $\pm$ 2.6 & 71.7 $\pm$ 2.2 \\
Shared & 16.3 $\pm$ 2.0 & 32.8 $\pm$ 3.9 & 49.5 $\pm$ 4.2 & 66.1 $\pm$ 3.4 \\
Attention (all agents) & 1.0 $\pm$ 0.2 & 3.0 $\pm$ 0.8 & 9.0 $\pm$ 3.1 & 15.3 $\pm$ 4.2 \\
Attention (2 nearest) & 1.0 $\pm$ 0.2 & 2.9 $\pm$ 1.1 & 8.3 $\pm$ 3.6 & 12.8 $\pm$ 5.7 \\
\midrule
\bottomrule
\end{tabular}
\end{table}

\begin{table}[t]
  \centering\small
  \caption{Minimum separation within an episode, averaged over episodes (higher is safer; the collision distance is
  $0.1$).}
  \label{tab:pt-minsep}
\begin{tabular}{l cccc}
\toprule
 & \multicolumn{4}{c}{Number of agents at test time} \\
\cmidrule(lr){2-5}
Policy (trained on 3 agents) & 3 (trained) & 4 & 5 & 6 \\
\midrule
Greedy (no avoidance) & 0.144 $\pm$ 0.000 & 0.089 $\pm$ 0.000 & 0.061 $\pm$ 0.000 & 0.045 $\pm$ 0.000 \\
\midrule
Independent & 0.192 $\pm$ 0.005 & 0.132 $\pm$ 0.003 & 0.095 $\pm$ 0.003 & 0.073 $\pm$ 0.004 \\
Shared & 0.201 $\pm$ 0.005 & 0.139 $\pm$ 0.005 & 0.104 $\pm$ 0.005 & 0.082 $\pm$ 0.004 \\
Attention (all agents) & 0.214 $\pm$ 0.008 & 0.172 $\pm$ 0.008 & 0.148 $\pm$ 0.006 & 0.131 $\pm$ 0.004 \\
Attention (2 nearest) & 0.214 $\pm$ 0.008 & 0.173 $\pm$ 0.008 & 0.150 $\pm$ 0.007 & 0.134 $\pm$ 0.007 \\
\midrule
\bottomrule
\end{tabular}
\end{table}

\begin{table}[t]
  \centering\small
  \caption{Percentage of time steps with a collision.}
  \label{tab:pt-violation}
\begin{tabular}{l cccc}
\toprule
 & \multicolumn{4}{c}{Number of agents at test time} \\
\cmidrule(lr){2-5}
Policy (trained on 3 agents) & 3 (trained) & 4 & 5 & 6 \\
\midrule
Greedy (no avoidance) & 2.80 $\pm$ 0.00 & 5.61 $\pm$ 0.00 & 8.76 $\pm$ 0.00 & 12.74 $\pm$ 0.00 \\
\midrule
Independent & 1.01 $\pm$ 0.11 & 2.07 $\pm$ 0.11 & 3.83 $\pm$ 0.25 & 5.64 $\pm$ 0.48 \\
Shared & 0.76 $\pm$ 0.11 & 1.71 $\pm$ 0.14 & 3.23 $\pm$ 0.09 & 4.96 $\pm$ 0.15 \\
Attention (all agents) & 0.04 $\pm$ 0.02 & 0.17 $\pm$ 0.09 & 0.42 $\pm$ 0.19 & 0.90 $\pm$ 0.31 \\
Attention (2 nearest) & 0.04 $\pm$ 0.02 & 0.16 $\pm$ 0.11 & 0.42 $\pm$ 0.23 & 0.95 $\pm$ 0.52 \\
\bottomrule
\end{tabular}
\end{table}

\begin{table}[t]
  \centering\small
  \caption{Success rate in percent (agents on their goal at the end of the episode).}
  \label{tab:pt-success}
\begin{tabular}{l cccc}
\toprule
 & \multicolumn{4}{c}{Number of agents at test time} \\
\cmidrule(lr){2-5}
Policy (trained on 3 agents) & 3 (trained) & 4 & 5 & 6 \\
\midrule
Greedy (no avoidance) & 100.0 $\pm$ 0.0 & 100.0 $\pm$ 0.0 & 100.0 $\pm$ 0.0 & 100.0 $\pm$ 0.0 \\
\midrule
Independent & 96.2 $\pm$ 0.6 & 96.3 $\pm$ 0.3 & 97.0 $\pm$ 0.2 & 97.1 $\pm$ 0.7 \\
Shared & 98.8 $\pm$ 0.5 & 98.9 $\pm$ 0.5 & 98.9 $\pm$ 0.6 & 98.8 $\pm$ 0.7 \\
Attention (all agents) & 97.4 $\pm$ 1.6 & 96.5 $\pm$ 1.6 & 94.9 $\pm$ 1.4 & 94.1 $\pm$ 1.3 \\
Attention (2 nearest) & 97.4 $\pm$ 1.6 & 96.2 $\pm$ 2.0 & 94.5 $\pm$ 1.9 & 92.9 $\pm$ 2.2 \\
\midrule
\bottomrule
\end{tabular}
\end{table}

\paragraph{Task success transfers, constraint satisfaction does not.}
Every learned policy reaches its goals in at least 92.9\,\% of the cases at every team size
(Table~\ref{tab:pt-success}), so the task itself transfers. The constraint does not. The share of episodes with a
collision (Table~\ref{tab:pt-collision}) is $16.3\,\%$ (shared) and $20.7\,\%$ (independent) with three agents and
roughly doubles when a fourth agent is added ($32.8\,\%$ and $38.7\,\%$); with six agents it is $66.1\,\%$ and
$71.7\,\%$. The mean minimum separation of the independent policies falls below the collision distance of $0.1$ at five
agents and that of the shared policy at six agents (Table~\ref{tab:pt-minsep}). A policy that looks successful by
its task reward can therefore already have lost much of its safety.

\paragraph{Part of the loss is density.}
The greedy controller, which has no avoidance at all, goes from $40.0\,\%$ to $92.8\,\%$ of episodes with a collision
between three and six agents, so a denser team is harder for any controller. The shared and independent policies
stay below the greedy controller at every size, but their advantage over it stays between $19$ and $34$ percentage points
and is smaller with six agents than with four: the increase from three to six agents is $49.8$ and $50.9$ points for
the shared and independent policies, against $52.8$ for the greedy controller. The advantage of the attention policy
over the greedy controller, in contrast, keeps growing, from $39$ to $78$ points.

\paragraph{The attention policy degrades much more slowly, and the cause is the architecture.}
The attention policy has $1.0\,\%$ of episodes with a collision at three agents and $15.3\,\%$ at six (an increase of
$14.3$ points), with a mean minimum separation of $0.131$ at six agents, still above $d_c$.
One might suspect that this is only because it sees all agents, whereas the other two see the two nearest. The
ablation rules this out for our setting: the same network fed only the two nearest neighbours reaches $12.8\,\%$ at
six agents (versus $15.3\,\%$; the difference is smaller than the seed spread of $4$--$6$ points) and is
indistinguishable at four and five agents (Table~\ref{tab:pt-collision}, Figure~\ref{fig:pt-transfer}). Seeing more
neighbours gave no measurable benefit. The comparison between the attention-on-two-nearest policy and the shared
policy differs in the architecture only, but ``architecture'' here is a bundle: set-structure with attention weights, an
explicit distance input $\|r_{ij}\|$, separate encoders, and about four times more parameters (Table~\ref{tab:pt-arch}).
We did not ablate these factors individually.

\paragraph{Shared versus independent.}
The shared policy is better than the independent policies in every column: $4.4$ to $8.9$ points fewer episodes with a
collision, a larger minimum separation and a higher success rate ($98.8\,\%$ versus $96.2\,\%$ at three agents). The
collision differences are of the same order as the seed spread ($1.4$--$4.2$ points), so we do not claim them to be
significant; the success difference is larger than the spread. Training curves
(Figure~\ref{fig:pt-train}) show that the shared policy learns faster than the independent ones, which is expected
since one network learns from the experience of all agents, and that the attention policy learns fastest.

\paragraph{Transfer gap.}
Table~\ref{tab:pt-gap} compares policies transferred from three to four agents with policies trained directly on four
agents. For the shared and independent policies the gap is large: $+17.7$ and $+12.9$ percentage points of episodes with a
collision. Directly trained policies are safer but succeed slightly less often, i.e.\ they trade some task success for
safety. For the attention policy there is no measurable gap ($-5.3$ points with a standard deviation of $8.7$); the
directly trained attention policies have $20.6\,\%$, $2.6\,\%$ and $1.8\,\%$ in the three seeds, so the poor average
comes from a single seed, and the transferred policy ($3.0\,\%$) is as safe as the two good ones. The reference runs use
the same training budget as the others and are therefore not necessarily converged.

\begin{table}[t]
  \centering\small
  \caption{Transfer gap at four agents: policies trained on three agents and transferred, versus policies trained
  directly on four agents. A positive gap means the transferred policy has more collisions (percentage points).}
  \label{tab:pt-gap}
\begin{tabular}{l cc c cc}
\toprule
 & \multicolumn{3}{c}{Episodes with a collision (\%)} & \multicolumn{2}{c}{Success rate (\%)} \\
\cmidrule(lr){2-4}\cmidrule(lr){5-6}
Policy & Transferred ($3\!\to\!4$) & Trained on 4 & Gap (pp) & Transferred & Trained on 4 \\
\midrule
Independent & 38.7 $\pm$ 1.4 & 25.8 $\pm$ 1.1 & +12.9 & 96.3 $\pm$ 0.3 & 92.0 $\pm$ 0.5 \\
Shared & 32.8 $\pm$ 3.9 & 15.1 $\pm$ 6.3 & +17.7 & 98.9 $\pm$ 0.5 & 96.9 $\pm$ 1.3 \\
Attention (all agents) & 3.0 $\pm$ 0.8 & 8.3 $\pm$ 8.7 & -5.3 & 96.5 $\pm$ 1.6 & 93.7 $\pm$ 4.0 \\
\bottomrule
\end{tabular}
\end{table}

\begin{figure}[t]
  \centering
  \includegraphics[width=0.92\linewidth]{figures/training_curves.png}
  \caption{Training curves on three agents (mean over three seeds, band = min to max; moving average of 10
  iterations). Left: training return. Right: fraction of time steps with a collision during training.}
  \label{fig:pt-train}
\end{figure}

\subsection{Discussion}
\label{sec:pt-discussion}
\paragraph{What this says about constraint-preserving transfer.}
Transfer of the task and transfer of the constraint are different things. The task survived a change of population
almost perfectly, while the constraint did not, and the loss was not predicted by task performance. For the
proposal this means that preservation of a constraint cannot be assumed from successful transfer; it has to be
measured directly. In our experiments, the set-based attention architecture preserved much more of the constraint
than the concatenated-observation policies, and its transfer gap at four agents was not measurable.

\paragraph{Information.}
When the team outgrows a $k$-nearest observation, each agent no longer observes every other agent. Our ablation
shows that, for the attention policy, observing all neighbours rather than only the two nearest did not improve the
measured collision outcomes. In this setting, the transfer gap is therefore not explained by this information
difference alone.

\paragraph{Practical reading.}
If a team may grow, the set-based attention policy is the most robust of the learned policies we tested: it has the
lowest collision rate while retaining high goal-reaching success across all evaluated team sizes.

\subsection{Limitations and threats to validity}
\label{sec:pt-limits}
\begin{itemize}
  \item \textbf{A small world.} The dynamics are discrete and noise free, there is no communication, and the arena is
        fixed, so the team becomes denser as it grows. Conclusions about larger teams, continuous dynamics or noise are
        not supported.
  \item \textbf{Three seeds.} We report means and standard deviations over three seeds without confidence intervals;
        only large effects (for example attention versus the other policy families) are reliable. The comparison between
        shared and independent policies is suggestive only.
  \item \textbf{Density and policy are confounded in $N$.} The greedy controller shows how large the density effect
        alone is; we did not run an experiment with an arena that grows with $N$.
  \item \textbf{``Architecture'' is a bundle.} The attention policy differs from the shared policy in set structure,
        attention weights, an explicit distance input and the number of parameters. We did not separate these factors,
        so we cannot say which one matters.
  \item \textbf{Reference runs.} The policies trained directly on four agents use the same 400 iterations as all others and
        one attention seed is clearly worse than the other two, so ``trained on 4'' is not a converged optimum.
  \item \textbf{Tuning.} The collision penalty was tuned on the same evaluation worlds with the shared policy
        (Section~\ref{sec:pt-training}); all other settings were fixed in advance.
  \item \textbf{Extension rule for independent policies.} Giving the new agent a copy of agent~1's network is one choice
        among several (random initialisation, an average of the trained networks); the results for independent
        policies depend on it.
  \item \textbf{Metrics.} Success is measured at the final step only, and evaluation uses greedy actions; stochastic
        execution was not studied.
\end{itemize}

\paragraph{Future work.}
(i) Train on a mixture of team sizes; (ii) isolate which components of the attention architecture matter;
(iii) test whether the findings persist with continuous dynamics, stochastic execution and larger arenas; and
(iv) combine the population axis with the spatial-scale axis studied in the other parts of the project.

\subsection{Reproducibility}
\label{sec:pt-repro}
All code is in the repository \texttt{marl-population-transfer} (files \texttt{env.py}, \texttt{policies.py},
\texttt{train.py}, \texttt{evaluate.py}, and \texttt{make\_figures.py}).
\texttt{run\_all.sh} trains all policies (three seeds, 400 iterations) and the four-agent references and produces
every table and figure of this section; it takes roughly 15--20 minutes on a laptop CPU. Evaluation uses a fixed
random seed (12345) so that all policies are tested on identical worlds, and every table entry in this section is
generated from the stored results file.


\bibliographystyle{plainnat}
\bibliography{refs_population_transfer}

\end{document}
